\subsection{Sprint 0}

    Para a sprint 0 a equipe definiu que faria toda a estrutura do projeto, deixando
Ages 1 primariamente com o Design, Ages 2 com o banco, Ages 3 com repositórios,
arquitetura, wiki e workflows e Ages 4 com a criação de um Board para atividades,
bem como a gestão da equipe e definição do escopo previsto para cada sprint. Da
minha parte me comprometi a ficar no Backend, teria então que definir tecnologia,
iniciar os repositórios, definir alguma arquitetura e convidar os colegas para os
repositórios.

    Olhando no panorama de equipe, a sprint 0 foi muito satisfatória, todas as
atividades previstas foram concluídas com sucesso, creio eu que fruto não apenas
de um bom trabalho de equipe, mas também de uma ótima gestão dos Ages 4, dava
para ver um engajamento geral da equipe. Da minha parte fiquei muito satisfeito
também com as atividades concluídas, utilizando uma tecnologia conhecida de
minha Ages 2, o FastAPI, fiz questão de proporcionar aos meus colegas uma
experiência bem melhor do que minha anterior, fiz questão de configurar um Docker
com compose e até a extensão de Dev Containers para garantir que não haja
complicações na hora de configurar o repositório.

    Além disso implementei um esqueleto de uma arquitetura híbrida, entre Clean e Camadas, para tentar aproveitar
o melhor dos dois mundos, trazendo a flexibilidade dos contratos com bancos de
dados da Clean, com uma complexidade de desenvolvimento semelhante à de
camadas, sempre levando em conta que estou definindo uma arquitetura que o
menos experiente dos programadores deve conseguir utilizar. Além disso adicionei
debugger para facilitar o ambiente de desenvolvimento, MakeFile para poder criar
alias, Alembic para migrations, SQLAlchemy para mapeamento de banco de dados e
entre outros. Por fim consegui implementar um workflow do GitHub, que roda os
testes ao abrir um pull request, a pedido de uma das Ages 4 que quer utilizar a
técnica de TDD. Por fim adicionei dois readmes, um bem detalhado explicando como
rodar o Backend, GitFlow e afins e o outro ensinando a adicionar uma feature do
zero no projeto.

    Depois de duas Ages, creio que posso dizer com confiança que esta foi uma
das melhores sprints que já vi decorrer. Espero de verdade que o time continue
assim, pois de momento não consigo identificar nenhum problema.

    Gosto muito de olhar logo acima, para pessoas que estão uma fase na frente
que eu fazendo um bom trabalho, como Ages 2 não tive tanto esta oportunidade e
portanto tive que aprender em outros lugares já que tinha 2 Ages 3 e eles não foram
tão participativos assim. Em compensação vejo que tenho muitas lições
interessantes a aprender de meus Ages 4 este semestre. Desde o primeiro dia já
chegaram com uma apresentação do termo de abertura pronto, além de algumas
propostas de como seguir implementando. Além disso, boa parte do tempo das
aulas está sendo ocupado em dinâmicas, onde o time discute muito o que entendeu
de termos de abertura e reunião com stakeholder, creio que estas conversas de
alinhamento facilitam o time a andar em sintonia.

    Além disso creio que a divisão dos Ages 4 para cada um auxiliar uma tarefa (design, banco, arquitetura e gestão) nesta
primeira sprint provou-se muito eficaz, pois apesar de ainda não ter muita conversa
entre a equipe, os Ages 4 funcionaram como uma ponte entre as equipes. Além
disso prevejo que esta proximidade de squads multinível que estamos prevendo,
causará uma aproximação entre integrantes de diferentes níveis de Ages que
melhorará ainda mais a comunicação da equipe.

    Vejo que com a arquitetura pronta, wiki e todos os workflow bem
implementados, fico muito dependendo de como a equipe anda, prevejo que terei de
dar uma boa ajuda aos colegas Ages 1 e 2 para entenderem a arquitetura e talvez
até alguma coisa de configuração de ambiente, que mesmo que tentei simplificar ao
máximo, talvez seja necessário auxiliar no download do Docker que é um pouco
mais complexo visto que no Windows precisa do WSL. A partir do momento que eu
conseguir fixar a ideia de arquitetura nos Ages 1 e 2, irei me voltar aos Ages 3 do
Frontend para entender como imaginam a integração.

    Num caso do Backend atrasar o Front, não tenho medo de ajudar um pouco com código, mas espero que esta
parte ande bem sozinha dado um tempo de aprendizado. No caso de tudo isso
funcionar, minha preocupação final seria relacionada ao deploy, gostaria de validar o
mais cedo possível que o Backend sobe na AWS e roda sem problemas, até para
eliminar um possível bloqueador que pode dar uma grande dor de cabeça, já que
por experiência própria, minhas 2 Ages passadas, falando com os Ages 3 perto do
final da Ages eles ainda estavam se preocupando com problemas no deploy, talvez
acabe entendendo eles, mas planejo tentar começar isso logo que possível para não
ficar com prazos apertados.